% ============================================================
% Body-text proof sketches (3-5 lines each) for T1--P7.
% Each theorem is stated in the body with these sketches;
% full proofs live in appendix_theorems.tex (\cref{app:proofs}).
% ============================================================

% --- T1 sketch (body) ---
% After stating Theorem T1:
\begin{proofsketch}
Relabeling is a bijection of corpora, so summing the equivariant
likelihood over the level set of a $G$-invariant statistic $T$ shows
$p(T\mid\sigma\cdot\theta)=p(T\mid\theta)$; with a $G$-invariant
prior, Bayes' rule transports this equality to the posterior, which is
therefore constant on orbits. Full proof in \cref{app:proofs}.
\end{proofsketch}

% --- Corollary 1.1 sketch (body) ---
% One line, inline:
% ``Since any measurable function of a $G$-invariant statistic is
% itself $G$-invariant, no downstream processing --- spectral,
% neural, or otherwise --- can break orbit-constancy (Cor.~1.1).''

% --- T2 sketch (body) ---
\begin{proofsketch}
By T1 the likelihood of any invariant statistic is constant on the
orbit, so $I(\theta;T)=0$ when $\theta$ is uniform on it; Fano's
inequality then lower-bounds the error of \emph{any} decoder by
$1-\log 2/\log N$. The finite-sample version replaces the exact orbit
by the $\varepsilon(n)$-empirical orbit whose radius comes from T3,
and the same argument prices the minimum external information:
$I(\theta;K)\ge(1-\delta)\log N(n)-\log 2$.
\end{proofsketch}

% --- T3 sketch (body) ---
\begin{proofsketch}
Lines are independent and each contributes a co-occurrence matrix of
spectral norm at most $2wL$, so the matrix Bernstein inequality
yields $\|\hat C-\mathbb{E}\hat C\|_2=O(wL\sqrt{\log d/m})$ with high
probability; the clipped PPMI map is entrywise Lipschitz, and Wedin's
$\sin\Theta$ theorem transfers the bound to the top-$k$ singular
subspaces whenever the spectral gap dominates. The bootstrap
stability we observe (CV $0.65\%$) and the failed sanity checks on
short-line corpora are the two regimes of this bound.
\end{proofsketch}

% --- T4 sketch (body) ---
\begin{proofsketch}
$B^{\top}A$ perturbs $(B^{\top}B)\Omega^{*}$ by $B^{\top}E$, and the
Procrustes solution is its orthogonal polar factor, whose
perturbation is controlled by $2\|B^{\top}E\|_F/\sigma_r(B)^2$
(Li 1995). A nearest-neighbour match is then correct for every test
pair whose margin exceeds
$2(\|x\|\|\hat\Omega-\Omega^{*}\|+\|e\|)$. Empirically the
instantiated threshold is vacuous on the synthetic benchmark
(certifies $0\%$ vs.\ observed $68.9\%$; margin AUC $0.536$) --- the
bound is sufficient, not necessary, and we report the gap
(App.~remark).
\end{proofsketch}

% --- L5 sketch (body) ---
% One-two lines, inline:
% ``The effective rank $e^{H(\sigma/\|\sigma\|_1)}$ inherits a
% Fannes--Audenaert modulus of continuity in the singular values
% (Lemma L5), so by Weyl's inequality it is robust to the sampling
% noise quantified in T3 --- our licence to use it as the primary C1
% statistic.''

% --- P6 sketch (body) ---
\begin{proofsketch}
A two-point exchange argument on the optimality conditions at
$\lambda_1<\lambda_2$ shows the decoded output's LM log-likelihood is
non-decreasing in $\lambda$; at $\lambda=1$ the decoder returns the
candidate-set maximizer, which is at least as probable as any real
sample in the pool. Hence a fusion decoder's bits-per-glyph sits
below the \emph{mean} cross-entropy of real lines by construction:
$7.31<8.26$ is a selection effect of the design, not evidence of
``super-realism''.
\end{proofsketch}

% --- D7/P7 sketch (body) ---
\begin{proofsketch}
C1, C2 and C3 are defined by what can falsify a statement: internal
null ensembles, nothing internal, or external anchors $K$. Since any
decision computed from invariant statistics is constant on orbits
(Cor.~1.1) while a C3 statement by definition varies within some
orbit, no anchor-free procedure decides a C3 statement
(Prop.~P7) --- the protocol is a theorem, not a style guide.
\end{proofsketch}

% Suggested environment definition (preamble):
% \newenvironment{proofsketch}
%   {\begin{proof}[Proof sketch]}{\end{proof}}
